% ============================================
% Johns Hopkins Letter Template
% Uses the built-in 'letter' class
% ============================================

\documentclass[11pt,letterpaper]{letter}

\usepackage{graphicx}
\usepackage{eso-pic}
\usepackage{geometry}
\usepackage{microtype}
\usepackage[utf8]{inputenc}
\usepackage[hidelinks]{hyperref}

% ----- Page Setup -----
\geometry{left=25mm,right=25mm,top=25mm,bottom=25mm}

% ----- Logo Placement (foreground, first page only) -----
\newcommand{\PutJHULogoFG}[1]{%
  \AddToShipoutPictureFG*{%
    \AtPageUpperLeft{%
      \hspace{10mm}%
      \raisebox{-38mm}{%
        \includegraphics[width=6cm]{#1}%
      }%
    }%
  }%
}

% ----- Sender Information -----
\signature{Decory Edwards}
\address{
Department of Economics \\
Johns Hopkins University \\
Baltimore, MD 21218 \\
\texttt{dedwar65@jh.edu}
}

\begin{document}
\PutJHULogoFG{Images/jhulogo.png}

\begin{letter}{
Search Committee \\
Department of Economics \\
Boise State University \\
Boise, ID
}

\opening{Dear Members of the Search Committee,}

I am writing to apply for the Visiting Assistant Professor position in the Department of Economics at Boise State University. I am a Ph.D.\ candidate in the Department of Economics at Johns Hopkins University (advisors: Christopher D.\ Carroll, Nicholas W.\ Papageorge, and Stelios Fourakis), and I expect to receive my degree in May 2026. My primary specializations are macroeconomic modeling, wealth and income dynamics, and computational economics.

My research agenda focuses on the ability of macroeconomic models to generate empirically realistic distributions of wealth and income over the life cycle. A key driver of that work is modeling heterogeneity in the rate of return on assets, and understanding how return dispersion contributes to portfolio accumulation across households.

In my job market paper, I calibrate a life cycle heterogeneous agent model to match wealth moments from the Survey of Consumer Finances by allowing households to differ in their portfolio returns. The model helps explain observed wealth concentration and return dispersion. In a related research project using the Health and Retirement Study, I examine how trust influences both income growth and asset returns. I find a hump shaped relationship for trust and returns (consistent with the literature) and characterize the distribution of the persistent component of returns to net worth. These experiences have equipped me to frame research strategies and design modeling approaches, execute econometric and simulation analysis with large micro data sets, and translate results into clear and rigorous written deliverables for both academic and practitioner audiences.

I am particularly attracted to Boise State’s emphasis on applied research across macroeconomics, public economics, labor, and growth. My work on taxation and heterogeneous returns connects directly to questions of economic growth, distribution, and policy design that align with the department’s stated fields. I maintain an active research agenda aimed at producing publishable work, and I would welcome the opportunity to contribute to a department that values both research productivity and high quality instruction.

Teaching is central to this position, and I am enthusiastic about contributing to a 3–3 teaching load. I have experience teaching and assisting in face to face university courses and am comfortable leading large sections of Principles of Economics. I emphasize clarity, structured problem solving, and real world applications when teaching introductory material. I would also welcome the opportunity to teach statistics or an upper level elective connected to macroeconomics, inequality, or public finance. The possibility of teaching multiple sections of principles courses aligns well with my experience and my commitment to undergraduate instruction.

I am enthusiastic about relocating to Boise and becoming part of the university’s academic community. The renewable structure of this appointment is especially appealing, as it would allow me to further develop my research agenda while contributing meaningfully to the department’s teaching mission over multiple years.

I believe I would be a strong fit for this position because:
\begin{enumerate}
    \item My research agenda aligns with the department’s strengths in applied macroeconomics and public economics.
    \item I have demonstrated experience teaching face to face courses and am committed to excellence in principles instruction.
    \item I am eager to contribute to a collaborative department that values both research and student success.
\end{enumerate}

Thank you very much for your consideration. I would welcome the opportunity to contribute to the Department of Economics at Boise State University.

\closing{Sincerely,}

\end{letter}
\end{document}
